\setpart{Qwen3-Omni · 输入与输出}

\begin{frame}[t]{Qwen3-Omni 有两种流式，要分开看}
\vspace{-.1cm}
\begin{minipage}[t]{5.0cm}\minipagetop
\tb{A. 模型内部的输出流（已有）}\\[.05cm]
\scalebox{.95}{\begin{tikzpicture}[x=1cm,y=1cm]
  \node[box=Teal,font=\scriptsize,minimum width=1.45cm,minimum height=.5cm] (t) at (.75,0) {Thinker};
  \node[box=Teal,font=\scriptsize,minimum width=1.45cm,minimum height=.5cm] (k) at (2.6,0) {Talker};
  \node[box=Teal,font=\scriptsize,minimum width=1.45cm,minimum height=.5cm] (c) at (4.45,0) {Code2Wav};
  \draw[strm] (t)--node[above,tiny]{token}(k);
  \draw[strm] (k)--node[above,tiny]{codec 块}(c);
\end{tikzpicture}}\\[.05cm]
\scriptsize 三个分区之间边生产边交接，输出音频分块返回。这是流水线，不是输入的更新。
\end{minipage}\hfill
\begin{minipage}[t]{5.0cm}\minipagetop
\tb{B. 模型外部的输入流（待研究）}\\[.05cm]
\scalebox{.95}{\begin{tikzpicture}[x=1cm,y=1cm]
  \node[box=Muted,font=\scriptsize,minimum width=1.2cm,minimum height=.42cm] (v) at (.6,.3) {视频帧};
  \node[box=Muted,font=\scriptsize,minimum width=1.2cm,minimum height=.42cm] (a) at (.6,-.25) {音频块};
  \node[box=Muted,font=\scriptsize,minimum width=1.2cm,minimum height=.42cm] (x) at (.6,-.8) {文本事件};
  \node[box=Gold,font=\scriptsize,minimum width=1.5cm,minimum height=.5cm] (b) at (2.7,-.25) {缓冲与对齐};
  \node[box=Teal,font=\scriptsize,minimum width=1.4cm,minimum height=.5cm] (th) at (4.55,-.25) {Thinker};
  \draw[flow] (v)--(b); \draw[flow] (a)--(b); \draw[flow] (x)--(b);
  \draw[flow] (b)--(th);
\end{tikzpicture}}\\[.05cm]
\scriptsize 持续到达的多模态输入，在推理期间如何进入同一次计算，是本汇报要回答的问题。
\end{minipage}
\par\takeaway{M* 的 Streaming\allowbreak{}GraphEdge 与 vLLM-Omni 的 async\_chunk 解决的是 A。B 是另一个问题：输入要进入一次正在进行的理解。}
\par\memo{来源：report v2 08a 章（三个分区与边的类型）；设计文档 §2.4（M* 的 StreamingGraphEdge 是分区之间的边）。B 的能力在三个框架中的状态见框架部分。}
\end{frame}

\begin{frame}[t]{Qwen3-Omni 的结构：输入经过编码器进入 Thinker}
\vspace{.05cm}
\begin{tikzpicture}[x=1cm,y=1cm,
  io/.style={box=Muted,font=\scriptsize,minimum width=2.2cm,minimum height=.42cm},
  enc/.style={box=Gold,font=\scriptsize,align=center,minimum width=2.4cm,minimum height=.6cm},
  lm/.style={box=Teal,font=\scriptsize,align=center,minimum height=.62cm},
]
  % 第一行：输入
  \node[io] (tx) at (1.3,0) {文本};
  \node[io] (au) at (4.9,0) {音频};
  \node[io] (im) at (8.5,0) {图像 / 视频};
  % 第二行：编码器（文本直接进入 Thinker）
  \node[enc] (ae) at (4.9,-1.35) {audio\_encoder\\AuT};
  \node[enc] (ve) at (8.5,-1.35) {vision\_encoder\\ViT};
  % 第三行：Thinker
  \node[lm,minimum width=9.0cm,minimum height=.9cm] (th) at (5.2,-2.95) {\textbf{Thinker}：MoE 混合专家 LLM · 30B-A3B · GPU 0--1（TP=2）\\文本、音频、图像 token 在同一个上下文里};
  % 第四行：Talker、Code2Wav、输出
  \node[lm,minimum width=2.9cm] (to) at (1.8,-4.25) {文本输出};
  \node[lm,minimum width=2.9cm] (tk) at (5.0,-4.25) {Talker（语音生成）\\3B-A0.3B · GPU 2};
  \node[lm,minimum width=2.6cm] (cw) at (8.5,-4.25) {Code2Wav（声码器）\\GPU 2};
  \node[io,minimum width=2.2cm] (ao) at (11.4,-4.25) {音频输出};
  % 连线
  \draw[flow] (tx.south) -- (tx.south |- th.north);
  \draw[flow] (au.south) -- (ae.north);
  \draw[flow] (im.south) -- (ve.north);
  \draw[flow] (ae.south) -- (ae.south |- th.north);
  \draw[flow] (ve.south) -- (ve.south |- th.north);
  \draw[flow] (th.south -| to.north) -- (to.north);
  \draw[flow] (th.south -| tk.north) -- (tk.north);
  \draw[strm] (tk.east) -- (cw.west);
  \draw[flow] (cw.east) -- (ao.west);
\end{tikzpicture}
\par\takeaway{输入先经过编码器变成 token，再与文本 token 一起进入同一个 Thinker。所以新的音频或视频不是独立请求，而是要进入 Thinker 的上下文。}
\par\memo{来源：基准仓库 report v2 08a 章（5880 上的 GPU 布局，不是本仓库的设定）。}
\end{frame}

\begin{frame}[t]{三类输入的到达节奏与时间语义}
\scriptsize
\renewcommand{\arraystretch}{1.15}
\begin{tabularx}{\textwidth}{>{\bfseries}p{2.3cm} p{2.6cm} p{2.4cm} Y}
\toprule
\rowcolor{Mist}
输入 & 到达节奏（设定的场景） & 数据形式 & 需要的时间语义 \\
\midrule
视频 & 30 帧/秒 & 单帧图像 & 帧的采集时间（PTS），用于视频的位置编码 \\
音频 & 每 20 毫秒一块（50 块/秒） & PCM 小块 & 每块覆盖的媒体时间区间（起点、终点） \\
文本：用户键入 & 事件，不固定 & 完整的一句 & 事件发生时间 \\
文本：ASR 转写 & 每 300--500 毫秒一段 & 文本片段 & 它对应的音频区间，不是转写完成的时间 \\
文本：系统上下文 & 事件 & 字幕、外部信息 & 生效的时刻 \\
\bottomrule
\end{tabularx}
\par\takeaway{同样是文本，三类的时间语义不同，不能放进同一个队列。ASR 的文本要按它对应的音频区间对齐。}
\par\memo{以上到达节奏是本汇报设定的测试场景，不是测量值。数值来源：设计场景中的 30 FPS 视频、20 毫秒音频块与 ASR 的延迟假设。}
\end{frame}

\begin{frame}[t]{时间窗口：对齐不等于同频率，也不等于每次都等齐}
\vspace{-.1cm}
\begin{tikzpicture}[x=1cm,y=.5cm]
  \def\X{1.9}
  \def\S{8.0}
  \foreach \tm/\lab in {0/0 ms,.5/500 ms,1/1000 ms}{
    \draw[Line,line width=.4pt] ({\X+\tm*\S},-3.0) -- ({\X+\tm*\S},.8);
    \node[font=\tiny,text=Muted] at ({\X+\tm*\S},-3.3) {\lab};
  }
  \node[anchor=west,font=\scriptsize\bfseries,text=Ink] at (0,.3) {视频 30 帧/秒};
  \foreach \i in {0,...,29}{\fill[Orange!70] ({\X+\i*\S/30-.04},.12) rectangle ({\X+\i*\S/30+.04},.42);}
  \node[anchor=west,font=\scriptsize\bfseries,text=Ink] at (0,-.5) {音频 20 毫秒块};
  \foreach \i in {0,...,49}{\fill[Teal!55] ({\X+\i*\S/50},-.7) rectangle ({\X+\i*\S/50+.09},-.4);}
  \node[anchor=west,font=\scriptsize\bfseries,text=Ink] at (0,-1.4) {ASR 文本段};
  \fill[Gold!70] ({\X+.60*\S},-1.6) rectangle ({\X+.88*\S},-1.3);
  \node[font=\tiny,text=Ink] at ({\X+.74*\S},-1.45) {900 ms 到达，对应 600--880 ms};
  \fill[Gold!25] ({\X+.45*\S},-1.6) rectangle ({\X+.6*\S},-1.3);
  \node[font=\tiny,text=Muted] at ({\X+.45*\S},-1.95) {600 ms};
  \draw[Red,line width=.9pt] ({\X+1.0*\S},.9) -- ({\X+1.0*\S},-3.0);
  \node[font=\tiny,text=Red,anchor=south] at ({\X+1.0*\S},.92) {t = 1000 ms：提交一次理解};
  \fill[Red!10] ({\X+0*\S},-2.3) rectangle ({\X+1.0*\S},-2.6);
  \node[font=\tiny,text=Red,anchor=west] at ({\X+.02*\S},-2.45) {窗口 [0, 1000 ms)};
\end{tikzpicture}
\par\takeaway{视频每 33 毫秒一帧、音频每 20 毫秒一块、文本按自己的时间段到达。窗口按媒体时间选取，每路取它覆盖的部分，不要求三路同时到达。}
\par\memo{示意图，时间为说明用，不是测量值。ASR 段的时间区间需要单独记录（音频区间与转写完成时间不同）。}
\end{frame}

\begin{frame}[t]{现有数据：Qwen3-Omni 单请求端到端（基准仓库）}
\scriptsize
\renewcommand{\arraystretch}{1.12}
\setlength{\tabcolsep}{3pt}
\begin{tabularx}{\textwidth}{>{\bfseries}p{3.4cm} Y Y Y}
\toprule
\rowcolor{Mist}
场景（5880，3 卡，公平档） & \textcolor{Teal}{M*} & \textcolor{Orange}{vLLM-Omni} & \textcolor{Blue}{SGLang-Omni} \\
\midrule
文本，1k 提示，256 token，端到端（毫秒，缓存关） & 1,645 & 1,613 & 1,596 \\
语音转文字，端到端（毫秒，缓存关） & 255 & 251 & 278 \\
语音转文字，端到端（毫秒，缓存开） & 252 & 665 & 308 \\
语音合成，每秒音频所需时间（毫秒） & 97 & 153 & 120 \\
\bottomrule
\end{tabularx}
\par\takeaway{这些是单个离线式请求的数据。它们说明三个框架在没有外部输入流时的起点，还不能说明输入更新的能力。}
\par\memo{来源：基准仓库 report v2 08b 章与 01 章，不在本仓库复测。文本与语音转文字为 11 月 11 日批次；缓存开的 M* 文本请求崩溃，未计入。}
\end{frame}
